% GENERATED FILE. Edit canonical lessons, then run npm run generate:books.
% canonical-source-hash: fnv1a64:06b9a445f884368d
% canonical-lessons: GE-C366-hypothese, GE-C366-theoretisch, GE-C366-unter-umstaenden, GE-C366-gegebenenfalls, GE-C366-wenigstens

\chapter{Theorie und Praxis}
\label{ch:ge-theorie-und-praxis}
\hlchaptermodality{366}

\begin{chapteropening}
\textbf{By the end of this chapter:} \emph{I can state a hypothesis, say what is true in theory, name the circumstances and say if applicable and at least.}

Complete the last lesson of chapter 366: I can state a hypothesis, say what is true in theory, name the circumstances and say if applicable and at least.
\end{chapteropening}

\section[die Hypothese]{die Hypothese — the hypothesis}

\label{lesson:GE-C366-hypothese}



\begin{quote}
Before the new one: say the German for apart from that, then the German for if need be.
\end{quote}

\subsection*{You'll want to know: die Hypothese}
\textbf{die Hypothese} — \textquotedblleft{}the hypothesis\textquotedblright{}.

\begin{quote}
Die Forscherin will ihre Hypothese mit einem Experiment prüfen.
\end{quote}

\begin{quote}
\emph{The researcher wants to test her hypothesis with an experiment.}
\end{quote}

Plural: \emph{die Hypothesen}. The stress is on the third syllable: \emph{hy-po-THE-se}.

\begin{cousinweb}
\emph{Hypothese} is a learned borrowing from Late Latin \emph{hypothesis}, from Greek \emph{hypóthesis} \textquotedblleft{}a base, the basis of an argument\textquotedblright{} (EtymDB), built from \emph{hypó} \textquotedblleft{}under\textquotedblright{} and \emph{thésis} \textquotedblleft{}a setting down\textquotedblright{}: what you put underneath.
\end{cousinweb}

\begin{grammarlens}[title={What you've built}]
One more word for telling, explaining and describing at B1.
\end{grammarlens}

\subsection*{Your turn}
\begin{itemize}
\raggedright
  \item \emph{Say it:} \emph{die Hypothese}
  \item \emph{Say it:} \emph{die Hypothese}, once more
  \item \emph{Say it:} a sentence of your own with \emph{die Hypothese}
  \item \emph{Recall:} say the German for apart from that, then the German for if need be, then say \emph{die Hypothese} again
\end{itemize}

\begin{tcolorbox}[breakable,colback=teal!4,colframe=teal!35!black,title={Before you move on}]
What does die Hypothese mean? (\textquotedblleft{}The hypothesis\textquotedblright{}.) Say it once more.
\end{tcolorbox}
\section[theoretisch]{theoretisch — in theory, theoretical}

\label{lesson:GE-C366-theoretisch}



\begin{quote}
Before the new one: say the German for if need be, then the German for the hypothesis.
\end{quote}

\subsection*{You'll want to know: theoretisch}
\textbf{theoretisch} — \textquotedblleft{}in theory, theoretical\textquotedblright{}.

\begin{quote}
Theoretisch kann man die Strecke in zwei Stunden fahren.
\end{quote}

\begin{quote}
\emph{In theory you can drive the route in two hours.}
\end{quote}

\emph{Theoretisch …, praktisch …} — in theory …, in practice …

\begin{cousinweb}
\emph{theoretisch} goes with \emph{Theorie}, a learned borrowing from Latin \emph{theōria} \textquotedblleft{}a philosophic speculation\textquotedblright{} (Lewis \& Short), from Greek \emph{theōría} \textquotedblleft{}a looking at\textquotedblright{} (EtymDB).
\end{cousinweb}

\begin{grammarlens}[title={What you've built}]
One more word for telling, explaining and describing at B1.
\end{grammarlens}

\subsection*{Your turn}
\begin{itemize}
\raggedright
  \item \emph{Say it:} \emph{theoretisch}
  \item \emph{Say it:} \emph{theoretisch}, once more
  \item \emph{Say it:} a sentence of your own with \emph{theoretisch}
  \item \emph{Recall:} say the German for if need be, then the German for the hypothesis, then say \emph{theoretisch} again
\end{itemize}

\begin{tcolorbox}[breakable,colback=teal!4,colframe=teal!35!black,title={Before you move on}]
What does theoretisch mean? (\textquotedblleft{}In theory, theoretical\textquotedblright{}.) Say it once more.
\end{tcolorbox}
\section[unter Umständen]{unter Umständen — under certain circumstances, possibly}

\label{lesson:GE-C366-unter-umstaenden}



\begin{quote}
Before the new one: say the German for the hypothesis, then the German for in theory.
\end{quote}

\subsection*{You'll want to know: unter Umständen}
\textbf{unter Umständen} — \textquotedblleft{}under certain circumstances, possibly\textquotedblright{}.

\begin{quote}
Unter Umständen dauert die Reparatur eine Woche.
\end{quote}

\begin{quote}
\emph{Under certain circumstances the repair may take a week.}
\end{quote}

Written short as \emph{u. U.} \emph{Der Umstand} — the circumstance; \emph{Umstände machen} — to cause trouble.

\begin{cousinweb}
\emph{Umstand} joins \emph{um} \textquotedblleft{}around\textquotedblright{} to \emph{Stand} \textquotedblleft{}standing, position\textquotedblright{} (EtymDB). It is usually explained as a rendering of Latin \emph{circumstantia} \textquotedblleft{}a standing around\textquotedblright{}, from \emph{circumstāre} \textquotedblleft{}to stand around\textquotedblright{} (Lewis \& Short).
\end{cousinweb}

\begin{grammarlens}[title={What you've built}]
One more word for telling, explaining and describing at B1.
\end{grammarlens}

\subsection*{Your turn}
\begin{itemize}
\raggedright
  \item \emph{Say it:} \emph{unter Umständen}
  \item \emph{Say it:} \emph{unter Umständen}, once more
  \item \emph{Say it:} a sentence of your own with \emph{unter Umständen}
  \item \emph{Recall:} say the German for the hypothesis, then the German for in theory, then say \emph{unter Umständen} again
\end{itemize}

\begin{tcolorbox}[breakable,colback=teal!4,colframe=teal!35!black,title={Before you move on}]
What does unter Umständen mean? (\textquotedblleft{}Under certain circumstances, possibly\textquotedblright{}.) Say it once more.
\end{tcolorbox}
\section[gegebenenfalls]{gegebenenfalls — if applicable, should the case arise}

\label{lesson:GE-C366-gegebenenfalls}



\begin{quote}
Before the new one: say the German for in theory, then the German for under certain circumstances.
\end{quote}

\subsection*{You'll want to know: gegebenenfalls}
\textbf{gegebenenfalls} — \textquotedblleft{}if applicable, should the case arise\textquotedblright{}.

\begin{quote}
Gegebenenfalls rufe ich Sie morgen noch einmal an.
\end{quote}

\begin{quote}
\emph{If necessary, I will call you again tomorrow.}
\end{quote}

Written short as \emph{ggf.}, very common on forms. Formal.

\begin{cousinweb}
\emph{gegebenenfalls} grew out of \emph{gegebenen Falls} \textquotedblleft{}in the given case\textquotedblright{}, from \emph{gegeben}, the participle of \emph{geben} \textquotedblleft{}to give\textquotedblright{}, and \emph{Fall} \textquotedblleft{}case\textquotedblright{}, Old High German \emph{val} (EtymDB), with the adverb ending \emph{-s}.
\end{cousinweb}

\begin{grammarlens}[title={What you've built}]
One more word for telling, explaining and describing at B1.
\end{grammarlens}

\subsection*{Your turn}
\begin{itemize}
\raggedright
  \item \emph{Say it:} \emph{gegebenenfalls}
  \item \emph{Say it:} \emph{gegebenenfalls}, once more
  \item \emph{Say it:} a sentence of your own with \emph{gegebenenfalls}
  \item \emph{Recall:} say the German for in theory, then the German for under certain circumstances, then say \emph{gegebenenfalls} again
\end{itemize}

\begin{tcolorbox}[breakable,colback=teal!4,colframe=teal!35!black,title={Before you move on}]
What does gegebenenfalls mean? (\textquotedblleft{}If applicable, should the case arise\textquotedblright{}.) Say it once more.
\end{tcolorbox}
\section[wenigstens]{wenigstens — at least}

\label{lesson:GE-C366-wenigstens}



\begin{quote}
Before the new one: say the German for under certain circumstances, then the German for if applicable.
\end{quote}

\subsection*{You'll want to know: wenigstens}
\textbf{wenigstens} — \textquotedblleft{}at least\textquotedblright{}.

\begin{quote}
Es hat geregnet, aber wenigstens war es warm.
\end{quote}

\begin{quote}
\emph{It rained, but at least it was warm.}
\end{quote}

Close to \emph{mindestens}, which mostly goes with numbers: \emph{mindestens drei Tage}.

\begin{cousinweb}
\emph{wenigstens} is built on the superlative of \emph{wenig} \textquotedblleft{}little\textquotedblright{}, from Middle High German \emph{weinec} (EtymDB), with the adverb ending \emph{-ens}: \textquotedblleft{}at the very least\textquotedblright{}.
\end{cousinweb}

\begin{grammarlens}[title={What you've built}]
One more word for telling, explaining and describing at B1.
\end{grammarlens}

\subsection*{Your turn}
\begin{itemize}
\raggedright
  \item \emph{Say it:} \emph{wenigstens}
  \item \emph{Say it:} \emph{wenigstens}, once more
  \item \emph{Say it:} a sentence of your own with \emph{wenigstens}
  \item \emph{Recall:} say the German for under certain circumstances, then the German for if applicable, then say \emph{wenigstens} again
\end{itemize}

\begin{tcolorbox}[breakable,colback=teal!4,colframe=teal!35!black,title={Before you move on}]
What does wenigstens mean? (\textquotedblleft{}At least\textquotedblright{}.) Say it once more.
\end{tcolorbox}
